\chapter{Tareas SQA}
\label{sec:tareas}

Este capítulo describe las tareas de aseguramiento de la calidad ejecutadas durante el sprint.

\section{Revisión de Productos de Software}
\label{sec:revision-productos}

SQA revisó los siguientes productos de software:

\begin{itemize}
    \item Código fuente del backend (Python/FastAPI) y frontend (TypeScript/Next.js).
    \item Documentación técnica (README, guías de despliegue, documentación de arquitectura).
    \item Transcripciones generadas (formato Markdown con metadatos YAML).
    \item Configuración de infraestructura (Docker, GitHub Actions, SonarCloud).
\end{itemize}

\section{Evaluación de Herramientas}
\label{sec:evaluacion-herramientas}

Se evaluaron las siguientes herramientas utilizadas en el proyecto:

\begin{longtable}{lcp{7cm}}
\caption{Evaluación de herramientas SQA}
\label{tab:herramientas-evaluadas}\\
\toprule
\textbf{Herramienta} & \textbf{Propósito} & \textbf{Resultado}\\
\midrule
\endfirsthead
\toprule
\textbf{Herramienta} & \textbf{Propósito} & \textbf{Resultado}\\
\midrule
\endhead
\bottomrule
\endfoot

Jira Cloud & Gestión de incidencias y bugs & Configurado con flujo To Do $\rightarrow$ In Progress $\rightarrow$ In Testing $\rightarrow$ Done.\\

SonarCloud & Análisis estático de código & Pipeline GitHub Actions configurado. Reporte inicial generado.\\

pytest & Pruebas unitarias e integración backend & Suite configurada en \texttt{backend/tests/}.\\

Cypress & Pruebas E2E frontend & Suite configurada en \texttt{web/tests/e2e/}.\\

Jest & Pruebas unitarias frontend & Suite configurada en \texttt{web/tests/unit/}.\\

\end{longtable}

\section{Verificación de Procesos}
\label{sec:verificacion-procesos}

SQA verificó los siguientes procesos mediante auditorías basadas en checklist:

\begin{enumerate}
    \item \textbf{Planificación y seguimiento:} Cumplimiento del sprint planificado (18-21 Julio).
    \item \textbf{Gestión de requisitos:} Trazabilidad entre requisitos funcionales y casos de prueba.
    \item \textbf{Diseño e implementación:} Calidad del código y aplicación de patrones de diseño.
    \item \textbf{Pruebas:} Cobertura de casos de prueba y reporte de defectos.
\end{enumerate}

\section{Auditorías de Configuración}
\label{sec:auditorias-configuracion}

Se verificó el control de configuración mediante:

\begin{itemize}
    \item Repositorio Git con rama dedicada \texttt{planning-SQA}.
    \item Flujo de trabajo basado en feature branches + pull requests.
    \item Pipeline de CI/CD en GitHub Actions para pruebas y análisis estático.
\end{itemize}

\section{Cronograma}
\label{sec:cronograma}

El sprint de calidad se ejecutó en 4 días según el siguiente cronograma:

\begin{longtable}{lcp{6cm}}
\caption{Cronograma del sprint SQA}
\label{tab:cronograma}\\
\toprule
\textbf{Día} & \textbf{Fecha} & \textbf{Actividades}\\
\midrule
\endfirsthead
\toprule
\textbf{Día} & \textbf{Fecha} & \textbf{Actividades}\\
\midrule
\endhead
\bottomrule
\endfoot

1 & 18 Jul & Planificación, configuración de Jira, análisis estático inicial (SonarCloud), diseño de casos de prueba y matriz de rastreabilidad.\\

2 & 19 Jul & Ejecución manual de pruebas, reporte de bugs, pruebas automatizadas backend (pytest) y frontend (Cypress/Jest).\\

3 & 20 Jul & Re-testing de bugs corregidos, análisis estático final (SonarCloud después), extracción de evidencias.\\

4 & 21 Jul & Consolidación del PDF final, grabación del video demo, cierre del informe.\\

\end{longtable}
